\section{Mức đáp ứng của hệ thống hiện tại}
\subsection{Cách đánh giá và phạm vi bằng chứng}
Đánh giá bằng đọc mã nguồn tại commit \texttt{cde7dcb}, ngày 07/09/2026; không suy ra trạng thái vận hành của cluster. Đơn vị đánh giá là \textbf{năng lực cần cho đề tài}, không dùng bộ thành phần hay thang trưởng thành của một khung kiến trúc MLOps.
\begin{itemize}
  \item \statusyes: Có implementation đáp ứng năng lực được mô tả ở hàng tương ứng; có thể kế thừa sau kiểm tra tích hợp.
  \item \statuspartial: Có phần nền nhưng còn thiếu dữ liệu, contract hoặc xử lý cần cho thí nghiệm.
  \item \statusno: Chưa tìm thấy implementation đáp ứng năng lực trong các luồng đã kiểm tra.
\end{itemize}
Mức “Đầy đủ” chỉ áp dụng cho yêu cầu hẹp trong bảng, không có nghĩa toàn bộ pipeline đã sẵn sàng hay được kiểm chứng production.

\begin{xltabular}{\linewidth}{@{}L{3.3cm} L{1.5cm} Y@{}}
\caption{Hiện trạng repo đối với đề tài đã chốt.}\label{tab:readiness}\\
\toprule
\textbf{Năng lực cần có} & \textbf{Mức độ} & \textbf{Bằng chứng hiện tại và phần cần bổ sung} \\
\midrule\endfirsthead
\toprule\textbf{Năng lực cần có} & \textbf{Mức độ} & \textbf{Bằng chứng hiện tại và phần cần bổ sung} \\
\midrule\endhead
Lưu prediction và identity & \statuspartial & E1--E2 có prediction ID, thời gian, tenant/project/version và features. Cần kiểm soát tính đầy đủ của batch; producer gửi sự kiện thành công theo cơ chế best-effort. \\
Xác suất lớp cho CBPE/ATC & \statuspartial & E3 chỉ trả max confidence thang 0--100. Cần xác suất lớp dương 0--1, class mapping, decision threshold và calibration version. \\
So sánh phân phối, xuất báo cáo & \statusyes & E4 chạy Evidently với reference/production, sinh báo cáo và summary. Kế thừa detector, pin phiên bản và hiệu chỉnh ngưỡng trên development stream. \\
Kiểm tra dữ liệu đầu vào & \statuspartial & E4 kiểm tra thiếu cột, null cao; gộp cột thiếu vào drift tổng hợp. Cần tách lỗi contract khỏi statistical drift và chặn batch không hợp lệ. \\
Cửa sổ dữ liệu tái lập & \statuspartial & E4 lấy các bản ghi mới nhất theo version. Cần chốt khoảng thời gian, danh sách prediction ID và snapshot reference trước chạy. \\
Kích hoạt tác vụ giám sát & \statuspartial & E5 có khóa, watermark theo số lượng bản ghi và idempotency cho DriftRun. Chưa có audit nhãn hay trigger retraining theo chất lượng. \\
Feedback và ngân sách nhãn & \statusno & E2/E5 chưa có nhận nhãn gắn prediction và quản lý mục đích, quota, nhãn đã tiết lộ. Cần LabelRequest, Feedback và budget ledger. \\
Ước lượng, kiểm chứng suy giảm & \statusno & Chưa có CBPE/ATC, baseline chất lượng theo version, khoảng bất định và kiểm chứng trên mẫu nhãn ngẫu nhiên. \\
Policy quyết định retraining & \statusno & Chưa có vòng KEEP, REQUEST\_LABELS, HOLD, RETRAIN với bằng chứng quyết định và giới hạn yêu cầu nhãn. \\
Chạy job với snapshot đầu vào & \statusyes & E6 tạo TrainingJob, snapshot source/data và submit qua execution backend. Tái sử dụng để chạy recipe cố định. \\
Tạo training set từ production & \statuspartial & E6 nhận file training data nhưng mặc định lấy workspace asset mới nhất. Cần nối nhãn và xuất snapshot từ pool production hợp lệ. \\
Đăng ký candidate sau build & \statusyes & E7 tạo ModelVersion và lưu artifact/metrics khi build thành công. Kế thừa lưu vết candidate. \\
Đánh giá candidate độc lập & \statusno & Metrics được ingest chưa thay thế evaluator tin cậy. Cần holdout riêng, so sánh champion và quyết định pass/fail/hold. \\
Benchmark ngân sách nhãn & \statusno & Chưa có replay che nhãn, policy đối chứng và thống kê chi phí/chất lượng của đề tài. \\
\bottomrule
\end{xltabular}

\subsection{Các điểm mã nguồn làm căn cứ}
Các đường dẫn tính từ gốc repo, chỉ dẫn tới chức năng đã kiểm tra:
\begin{itemize}
  \item \textbf{E1:} \codepath{services/model-server/src/index.py}, gửi event và endpoint predict.
  \item \textbf{E2:} \codepath{services/consumer/src/database.py}, production logs và automatic-drift outbox.
  \item \textbf{E3:} \codepath{services/machine-learning-serving/src/index.py}, xử lý \codepath{predict\_proba}.
  \item \textbf{E4:} \codepath{services/evidently/src/main.py}, \codepath{load\_production\_data} và \codepath{run\_drift\_analysis}.
  \item \textbf{E5:} \codepath{services/control-plane/src/apps/drift/services/automatic.py}, tạo DriftRun.
  \item \textbf{E6:} \codepath{services/control-plane/src/apps/training/services/jobs.py}, tạo/submit job và snapshot.
  \item \textbf{E7:} \codepath{services/control-plane/src/apps/registry/services/versions.py}, đăng ký version.
\end{itemize}
Nền tảng có các khối suy luận, drift report, chạy training và lưu candidate. Khối cần xây nằm ở \textbf{feedback, ngân sách, kiểm chứng và quyết định}; không cần phát triển thêm hạ tầng AWS/Kubernetes để bắt đầu thực nghiệm.
